\section{18. Cluster Stock, Conservative Contact, and Lensing}
\label{sec:clusterstock}

Cluster lensing and dynamics probe the same gravitating inventory. In the
weak-field kinetic realization of O.3, baryons and a positive retained stock
move in a common potential,
\[
\nabla^2\Phi=4\pi G(\rho_*+\rho_g+\rho_C),\qquad
\kappa=\frac{\Sigma_*+\Sigma_g+\Sigma_C}{\Sigma_{\rm crit}}.
\]
The static response prescribes a possible equilibrium of this stock.
Adding that equilibrium as a second halo would duplicate its mass.
Gas can shock while stars and weakly coupled stock continue through an
encounter. A spatial separation of gas and lensing peaks is therefore
dynamically possible; its size depends on the initial distributions,
encounter, and contact law.

\paragraph{One cohort budget.}
For a fully released closed cohort, use the conservative normalization
\[
k_C=\frac{e^{-\eta_*}}{\epsilon}=5.1760752,\qquad C_0=k_C B_0,
\qquad f_{b,\rm coext}=\frac1{1+k_C}=0.1619151.
\]
Here $B_0$ is its ancestral baryonic mass and $C_0$ its surviving stock.
The mass convention and release history must be kept fixed across systems.
Inside an observed aperture, the retained fractions of baryons and stock
can differ. Writing them as $f_b^{\rm ret}$ and $f_C^{\rm ret}$ gives
\[
\frac{C}{B}=k_C\frac{f_C^{\rm ret}}{f_b^{\rm ret}},\qquad
\eta_{\rm ret}:=\frac{k_C B}{M_{\rm grav}-B}
=\frac{f_b^{\rm ret}}{f_C^{\rm ret}}.
\]
Thus $M_{\rm grav}=(1+k_C)B$ tests coextensiveness. It is a restrictive
special case of the transport model. For any positive inferred stock
$C=M_{\rm grav}-B$, the minimal compatible closed inventory is
\[
B_0=\max(B,C/k_C),\qquad
B_{\rm out}=B_0-B\ge0,\quad C_{\rm out}=k_CB_0-C\ge0.
\]
This construction supplies an admissible inventory. Its assembly history,
work cost, and transport time require a forward solution.

\paragraph{Reciprocal contact with finite available work.}
The coupling of O.3 admits central pair transitions in the relative
radial velocity $w$. For a pair with reduced mass
$\mu=m_im_j/(m_i+m_j)$ and separation direction $\boldsymbol n$,
\[
\delta\boldsymbol v_i=\frac{m_j}{m_i+m_j}\delta w\boldsymbol n,
\qquad
\delta\boldsymbol v_j=-\frac{m_i}{m_i+m_j}\delta w\boldsymbol n,
\qquad
\Delta K=\frac\mu2[(w+\delta w)^2-w^2].
\]
Both linear and angular momentum are conserved by each impulse. The work
and heat registers balance $\Delta K$ event by event. A symmetric factor
$v(U)v(U-W)$ suppresses a transition when either endpoint has exhausted
its available work; the floor state permits no such transition. Local
partners include gas and stars from both progenitors. An initial cohort
label does not confer a special interaction eligibility.

Seventeen controlled runs comprise fourteen encounters and three isolated
controls. They use coextensive Plummer progenitors of mass ratio $0.3$,
the same $k_C$, a gas adiabatic index $5/3$, and a selected two-component
relative-energy target. Weak contact preserves a stock concentration near
the stars while the gas lags. Increasing the rate alters that separation,
but the strongest tested contact also changes the isolated stock median
radius by $-27\%$ and the gas median radius by $+46\%$ within six dynamical
time units. It has not selected a stable cluster equilibrium.
The zero-work run permits zero events and reproduces the no-contact
trajectory; a low-work run permits 61 events against 7114 with the larger
work reserve. The largest energy drift is below $0.007\%$ of the binding
energy, with momentum errors at roundoff. Resolution, timestep, and
contact-range controls support the qualitative mechanism. These are
selected effective encounters; their rates, target, and work reserves
have not been inferred independently from cluster data. The spatial
stress of the full depleted scalar is not evolved in this particle test.

\paragraph{Radial data and independent lensing.}
Public X-COP profiles provide twelve radial targets, seven with stellar
profiles; five overlap the MENeaCS weak-lensing sample
\citep{Ettori2019Cluster,Herbonnet2020Cluster}. With one fixed normalization,
all nominal NFW targets admit positive cumulative stock and positive
stock in the fifteen tested radial shells. The forward hydrostatic
reconstruction of A644 has negative inferred shells, so positivity depends
on a physically admissible mass reconstruction. Projected stellar masses
are approximately deprojected; unmeasured diffuse light and correlated
profile errors remain relevant.

For the five independent lensing targets, the self-consistent coextensive
$M_{500}$ ranges, in $10^{14}M_\odot$, are
\[
\begin{array}{c|cc}
 &M_{500}^{\rm coext}&M_{500}^{\rm WL}\\\hline
\mathrm{A1795}&3.96\text{--}4.23&9.3\pm2.2\\
\mathrm{A2029}&7.46\text{--}8.07&12.1\pm2.5\\
\mathrm{A2142}&9.29\text{--}10.02&9.7\pm2.3\\
\mathrm{A85}&5.20\text{--}5.43&5.7\pm2.2\\
\mathrm{ZW1215}&4.19\text{--}4.46&3.5\pm2.2
\end{array}
\]
The ranges double the nominal stellar mass and are sensitivity ranges.
A1795 and A2029 require substantial relative baryon--stock segregation
under the central lensing targets: $\eta_{\rm ret}=0.518$ and $0.704$.
A positive transport inventory can accommodate these values, but it does
not predict them. ZW1215 changes substantially with the lensing estimator;
the two estimators use overlapping data and cannot be counted as independent
confirmations. The empirical companion gives the aperture conventions
and profile-level qualifications.

\paragraph{Gas support and possible outlier mechanisms.}
Retaining the common potential fixes the local gas balance,
\[
a_g=-g-\frac{P_{\rm th}'}{\rho_g}
       -\frac{P_{\rm nt}'}{\rho_g}+\frac{f_{Cg}}{\rho_g},\qquad
M=M_{\rm th}+M_{\rm nt}+M_{\rm contact}-M_{\rm acc},
\]
where $M_{\rm contact}=r^2f_{Cg}/(G\rho_g)$ and
$M_{\rm acc}=r^2a_g/G$. Assembly segregation, acceleration, and reciprocal
contact are possible contributions. A quiet gas velocity field alone
does not determine all of these terms. A passive drag vanishes at zero
relative drift, however, and a persistent supporting force requires its
own constitutive derivation. The small measured turbulent support in
A1795 and A2029 constrains a purely turbulent explanation
\citep{Sarkar2026Cluster,XRISM2025Cluster}. These possibilities remain
conditional pending joint lensing, pressure, and transport fits.

\paragraph{A resolved merger target.}
The public CANUCS Bullet v1 reconstruction supplies a separate spatial
test \citep{Rihtarsic2026Cluster}. At $z_l=0.296$, in its
$H_0=70\,\mathrm{km\,s^{-1}\,Mpc^{-1}}$, $\Omega_m=0.3$ convention,
$\Sigma_{\rm crit}=1.8268761\times10^9M_\odot\,\mathrm{kpc}^{-2}$
for $D_{ls}/D_s=1$. Within 150 kpc of the Bullet's BCG3, the reconstructed
total mass is $8.511\times10^{13}M_\odot$ and the supplied gas mass is
$6.184\times10^{12}M_\odot$. The released subcluster halo centre lies
3.92 kpc from BCG3, while the supplied gas-density peak lies about
169 kpc away. The latter is a pixel-based gas-density diagnostic.
The non-gas map still includes stars. These are targets extracted from
an existing best-fit reconstruction, with no simulated-map likelihood.
The illustrative encounters have 24--30 kpc force softening and roughly
90 kpc map smoothing under the quoted physical scaling; they cannot test
the observed few-kiloparsec alignment. A common selected contact must
reproduce the isolated profiles and the merger morphology together.

